\chapter{The Model}\label{ch:model}

We consider a finite mixture of Gaussian components, which is the model whose
estimation problem is studied in this report. The background facts from
Chapter~\ref{ch:background} are used throughout.

\section{Mixture density}\label{sec:mixture}

Given weights $\pi_1,\dots,\pi_k$ with $\pi_j\ge 0$ and $\sum_j \pi_j = 1$, the
density of an observation is
\begin{equation}\label{eq:mixture}
  p(x \mid \theta) = \sum_{j=1}^{k} \pi_j\, \phi\bigl(x;\,\mu_j,\Sigma_j\bigr),
\end{equation}
where $\phi(\cdot;\mu,\Sigma)$ denotes the multivariate normal density and
$\theta = \{\pi_j,\mu_j,\Sigma_j\}_{j=1}^{k}$ collects all parameters.

\lipsum[11-12]

\section{Identifiability}\label{sec:identifiability}

The mixture in Equation~\eqref{eq:mixture} is identifiable up to relabelling of
the components, provided the component means are distinct and the mixing
weights are strictly positive. Section~\ref{sec:prob} supplies the probabilistic
language needed for this statement.

\lipsum[13-14]

\section{Likelihood}\label{sec:likelihood}

For an independent sample $x_1,\dots,x_n$ the log-likelihood is
\begin{equation}\label{eq:loglik}
  \ell(\theta) = \sum_{i=1}^{n} \log \sum_{j=1}^{k} \pi_j\,\phi(x_i;\mu_j,\Sigma_j).
\end{equation}
Maximising Equation~\eqref{eq:loglik} directly is not convenient because the
logarithm does not act on individual terms. The estimation procedure of
Chapter~\ref{ch:estimation} avoids this difficulty.

\lipsum[15-16]

\begin{table}[ht]
  \centering
  \begin{tabular}{lcc}
    \toprule
    Component & Parameters & Free scalars \\
    \midrule
    Weight & $\pi_j$ & $1$ \\
    Mean & $\mu_j$ & $d$ \\
    Covariance & $\Sigma_j$ & $d(d+1)/2$ \\
    \bottomrule
  \end{tabular}
  \caption{Free parameters contributed by each component.}\label{tab:params}
\end{table}

Table~\ref{tab:params} counts the free parameters, which determine the cost
of the estimation problem discussed in Section~\ref{sec:counting}.

\section{Parameter counting}\label{sec:counting}

Summing the entries of Table~\ref{tab:params} over $k$ components, and
subtracting one for the normalisation constraint on the weights, gives
\begin{equation}\label{eq:count}
  N(k,d) = k\Bigl(1 + d + \tfrac{d(d+1)}{2}\Bigr) - 1 .
\end{equation}
Equation~\eqref{eq:count} grows quadratically in $d$ and linearly in $k$.

\lipsum[17-18]

\section{Further remarks}\label{sec:further-model}

The remarks in this section restate the main point of the chapter in more
general terms, so that the argument can be read without the surrounding
derivations.

\lipsum[75-86]
